Evaluating generated artifacts requires more than assigning a score: an evaluator
must obtain relevant evidence, apply task-specific criteria, and reconcile
competing aspects of quality. A disagreement with human preference may arise at
any of these stages, making the evaluator itself a useful object of optimization.
We study learning executable evaluation programs from human overall preferences
while keeping the underlying model fixed. We introduce \ERA{}, Evidence-Guided
Revision of Agentic Evaluators, as the method for revising these programs.
An evaluation program comprises artifact representations, rubrics,
specialist evaluators, evidence tools, routing, and score synthesis. \ERA{} uses
artifact-level execution evidence to formulate revision hypotheses, constructs
bounded program edits, and assesses candidates on a separate training cohort.
At deployment, the resulting evaluator scores each artifact independently,
without access to its opponent or preference labels. We specify an evaluation
protocol that separates adaptive development from terminal validation and held-out
testing, and tests both preference agreement and the usefulness of evaluator
feedback for artifact refinement. The proposed experiments assess preference
agreement, fixed-pool selection, and feedback-guided refinement under matched
models and budgets.
